\section{Related Work}
\label{sec:related-work}

\subsection{ConvergenceSweep (B.16)}

\cs{}~\citep{b16paper} introduced the seed-walk approach to minimum-edge-cut
redistricting.
Starting from the content-derived seed $s_0 = \seedfmt{\texttt{census\_release\_id}
\,\|\, \texttt{"DIA\_SEED\_V1"}} \bmod 2^{31}$, \cs{} advances one seed at
a time, tracking the normalised edge cut $\ECnorm$ achieved by each METIS
call.
The sweep halts after $T$ consecutive seeds produce no improvement.
The final plan is the one achieving the global minimum $\ECnorm$ across all
seeds tried.

B.16's central contribution is empirical certification that $T = 600$ is
sufficient for all 50 states.
Georgia is the hardest case: the algorithm last improves at seed $s_0 + 489$,
then runs 511 consecutive non-improving seeds before halting.
A threshold of $T = 500$ would fail Georgia.
The recommended $T_{\mathrm{stat}} = 600$ provides an 89-seed margin with
$P(\mathrm{tail} > 600) < 0.001$ under a fitted Gumbel extreme-value model.

\ps{} is a direct generalisation of \cs{}.
Where \cs{} always returns the rank-0 plan (minimum edge cut), \ps{}
parametrises the return rank.
The $T$-seed infrastructure is identical; only the selection rule changes.

\subsection{Seed Sensitivity (B.7)}

B.7~\citep{b7paper} characterised the full solution space of METIS
redistricting through a 50-state, 10,000-seed sweep (500,000 total METIS
calls).
The key findings for \ps{} are:
\begin{enumerate}
\item The coefficient of variation of $\ECnorm$ across seeds is below 2\%
      for 48 of 50 states.
      Georgia and North Carolina exhibit higher variance (CV $= 4.3\%$ and
      $3.8\%$ respectively) due to non-power-of-two district counts and
      geographic complexity.
\item Partisan seat-share variance across 10,000 seeds is at most 2 seats
      for the two highest-variance states.
      For most states, variance is zero or one seat.
\item The DIA single-seed specification is within $2.9\% \pm 1.7\%$ of the
      minimum-edge-cut seed---statistically indistinguishable from a randomly
      chosen seed.
\end{enumerate}

Finding (2) is the direct precursor to H.0's insensitivity result.
If varying the seed over 10,000 draws produces at most 2 seats of partisan
variation, varying the rank selection percentile over the sorted distribution
of 101 draws should produce even less variation.

\subsection{GerryChain Ensemble Positions (G.1)}

G.1~\citep{g1paper} placed the \ar{} + \cs{} plan within real \gc{} \recom{}
ensembles for six states.
The comparison metric is normalised edge-cut fraction.
The key findings:
\begin{itemize}
\item Wisconsin: 0.2nd percentile of the ensemble (more compact than 99.8\%
      of valid plans).
\item Georgia: 0.1st percentile.
\item Pennsylvania: 0.2nd percentile.
\item North Carolina: 50th percentile.
      Geographic structure constrains the feasible space so tightly that
      the minimum-cut solution and the ensemble median are nearly identical
      (plan cut $= 0.0973$, ensemble mean $= 0.0970$, $\Delta < 0.3\%$).
\item Texas and California: both expected below the 5th percentile.
\end{itemize}

These positions establish the baseline for the over-optimising objection and
motivate \ps{}'s percentile parametrisation.
For states at the 0.1st percentile, the legal question is whether being more
compact than 99.9\% of plans is an asset or a liability.
\ps{} makes this a statutory choice.

\subsection{Parameter Sensitivity (B.17)}

B.17~\citep{b17paper} conducted a 50-state sweep over five algorithm
parameters and found that partisan seat counts are insensitive to parameter
tuning within a fixed algorithm family.
The maximum variation across the full parameter range is $\Delta
D\text{-seats} \leq 0.3$ nationally.
The key parameter relevant to H.0 is the \cs{} threshold $T$: B.17 found
no measurable effect on partisan outcomes for $T \in [100, 1{,}000]$.

H.0 extends this result.
If varying $T$ from 100 to 1,000 does not shift partisan outcomes, then
varying the rank selected from within the resulting distribution should
similarly leave partisan outcomes unchanged.
The empirical results in Section~\ref{sec:results} confirm this.
